\section{Schemat ideowy i integracja}
\lead{Wspólny widok zdarzeń przy zachowaniu odrębności funkcji pożarowych.}
\fig{architecture.pdf}{1}{Architektura funkcjonalna. Niebieski: sieć zabezpieczeń; czerwony: autonomiczny SSP; linia przerywana: eksport informacji.}
\subsection*{Co dzieje się po alarmie?}
Zdarzenie KD lub SSWiN trafia do platformy integracyjnej w SOC, która pokazuje miejsce, czas, status urządzenia i przypisaną kamerę. Operator potwierdza obsługę i prowadzi kartę incydentu. Klip dowodowy eksportuje osoba z odpowiednią rolą, a operacja pozostawia ślad w dzienniku.

SSP wykonuje lokalny scenariusz niezależnie od platformy integracyjnej. Bramka SSP udostępnia zdarzenia przez funkcję tylko do odczytu, zatwierdzoną przez producenta. W konfiguracji SOC nie udostępnia się kasowania alarmów, blokowania czujek ani zmiany sterowań pożarowych.

\begin{note}{Praca przy awarii}
\textbf{Brak LAN:} KD używa lokalnych uprawnień; SSWiN działa lokalnie i zgłasza zdarzenia kanałem zapasowym; kamery zapisują na kartach. \textbf{Brak VMS/AI:} utrata wspólnego widoku, lecz nie lokalnej detekcji SSP. \textbf{Uszkodzenie centrali:} dotknięty budynek wymaga procedury zastępczej i interwencji serwisu.
\end{note}
